\centerline{\bf \Huge Взвешивания}

\begin{flushright}
        \scriptsize{Из любой, даже самой плохой ситуации бывает выход.}
\end{flushright}

\problem{0.1} 
Есть 3 одинаковых монет, из которых одна - фальшифая (более лёгкая). Как с помощью чашечных весов без гирь выделить фальшивую монету за наименьшее число взвешиваний.

\problem{0.2}Среди четырёх монет имеется одна фальшивая, отличающаяся от всех остальных по весу. Как выделить её спомощью двух взвешиваний на чашечных весах без гирь? Можно ли при этом выяснить, легче ли она, чем остальные монеты?

\problem{1} Из 9 монет одна — фальшивая, она тяжелее настоящих. Как найти ее за два взвешивания?

\problem{2} Из 27 монет одна — фальшивая, она легче настоящих. Найдите ее за 3 взвешивания.


\problem{3} Среди 101 монеты есть одна фальшивая, которая по весу отличается от настоящей. Но на этот раз неизвестно, в какую сторону. За два взвешивания определите, легче или тяжелее настоящей фальшивая монета. (Саму монету определять не нужно.)


\problem{4}
Есть 9 кг крупы и чашечные весы с гирями в 50г и 200 г.Как в три приема отвесить 2 кг крупы?


\problem{5}  
Из 21 монеты 10 настоящих и 11 фальшивых, причём каждая фальшивая на 1 г легче каждой настоящей. Взяли одну из монет. Как за одно взвешивание на весах со стрелкой определить фальшивая ли монета?

